\documentclass[11pt,a4paper,landscape]{article}
\usepackage[T1]{fontenc}\usepackage[utf8]{inputenc}\usepackage[french]{babel}
\usepackage{lmodern,amsmath,amssymb,graphicx,booktabs,xcolor,hyperref,fancyhdr}
\usepackage[margin=16mm]{geometry}
\hypersetup{colorlinks=true,linkcolor=blue,urlcolor=blue,pdftitle={Reprise affine puis RBF contrainte par les gradients}}
\graphicspath{{../results/hudge_rbf_restart_20261115/figures/}}
\pagestyle{fancy}\fancyhf{}\fancyhead[L]{WGPU-Matching — Reprise affine puis champ non affine}
\fancyhead[R]{15 novembre 2026}\fancyfoot[R]{\thepage}\setlength{\headheight}{14pt}
\setlength{\parindent}{0pt}\setlength{\parskip}{6pt}
\newcommand{\fig}[2]{\begin{center}\includegraphics[width=\linewidth,height=#2,keepaspectratio]{#1}\end{center}}
\begin{document}
{\LARGE\bfseries Affine puis fonctions de base radiale :\par
apport des gradients et du contrôle de cohérence\par}
\vspace{4mm}
\textbf{Expérience exécutée :} 191 paires indépendantes référence/cible de la plaque
trouée, trois représentations du champ, trois seuils de cohérence, soit 1\,719 configurations.
Chaque paire comporte 36 fenêtres de $32\times32$ pixels. Les cibles successives ne sont
jamais utilisées comme références intermédiaires.

Les fonctions de base radiale (RBF) donnent effectivement une meilleure représentation
non affine des observations disponibles. Utiliser les gradients des fenêtres comme
contraintes souples améliore aussi la prédiction des observations laissées hors ajustement.
\textbf{Mais ce gain ne crée pas de fenêtre retenue supplémentaire par rapport aux RBF
ajustées aux seuls déplacements dans cette campagne.} Le seuil de cohérence du gradient
est ici le facteur dominant.

\begin{center}\begin{tabular}{lrrr}
\toprule
Champ proposé après la phase affine & Seuil 0,25 & Seuil 0,50 & Seuil 1,00\\\midrule
Nouvelle affine, essais sur tous les centres &6\,035&6\,549&6\,762\\
RBF sur déplacements &6\,056&6\,549&6\,772\\
RBF sur déplacements et gradients &6\,056&6\,549&6\,772\\\bottomrule
\end{tabular}\end{center}
Nombres de fenêtres \emph{retenues}, sur 6\,876. La phase affine archivée commune en
retient 5\,804. Les différences au seuil strict sont donc : +231 par la nouvelle reprise
affine, puis +21 supplémentaires par les RBF. Ce ne sont pas 252 gains propres aux RBF.

\textbf{Au niveau 191 :} 18 fenêtres au départ, puis 25 avec la nouvelle affine au seuil
strict, 27 avec les RBF ; les trois représentations atteignent 31 au seuil 0,50 et 34 au
seuil 1,00. Les deux dernières fenêtres restent refusées par le solveur.

\textbf{Limite scientifique :} les fichiers de déplacements disponibles de l'éprouvette
en croix ne correspondent pas à cette plaque. Ces résultats mesurent l'acceptation
numérique et photométrique, pas une exactitude mécanique vérifiée. Les 191 niveaux
d'une même plaque ne constituent pas 191 réalisations indépendantes de matériau.

\newpage
\section*{1. Un seul champ dont on dérive les gradients}
Soient les centres validés $\mathbf x_i$, leurs déplacements $\mathbf d_i$ et leurs
gradients $G_i=F_i-I$. L'affine globale $A\mathbf x+\mathbf t$ est ajustée de façon robuste
aux positions cibles $\mathbf x_i+\mathbf d_i$. On la complète par des gaussiennes :
\[
 \widehat{\mathbf u}(\mathbf x)=(A-I)\mathbf x+\mathbf t+
 \sum_{j=1}^{n}\mathbf c_j\,b_j(\mathbf x),\qquad
 b_j(\mathbf x)=\exp\!\left[-\frac{\|\mathbf x-\mathbf x_j\|^2}{2\ell^2}\right],\quad \ell=64\ \mathrm{pixels}.
\]
La jacobienne est calculée analytiquement à partir de cette même fonction :
\[
 \widehat G(\mathbf x)=A-I-\sum_j\mathbf c_j\otimes
 \frac{\mathbf x-\mathbf x_j}{\ell^2}\,b_j(\mathbf x),\qquad
 \widehat F=I+\widehat G.
\]
Les composantes du gradient ne sont donc pas interpolées indépendamment du déplacement.
Pour ajuster les coefficients $\mathbf c_j$, on minimise :
\[
 \sum_i\|\widehat{\mathbf u}(\mathbf x_i)-\mathbf d_i\|_2^2
 +L_g^2\sum_i\|\widehat G(\mathbf x_i)-G_i\|_F^2
 +\lambda\sum_j\|\mathbf c_j\|_2^2,
 \qquad L_g=16\ \mathrm{pixels},\quad\lambda=10^{-4}.
\]
Le terme de gradient est absent dans la variante RBF déplacements. $L_g$ donne aux
écarts de gradient une échelle de déplacement sur un demi-support. Les contraintes
sont \textbf{souples}, car les gradients résolus ne sont pas exempts d'erreur. La
pénalisation des coefficients stabilise le système ; l'affine est fixée avant l'ajustement
du résidu RBF. Les valeurs sont fixées avant la campagne, sans optimisation sur les résultats.

\textbf{Boucle.} Ajuster sur les seules fenêtres retenues ; proposer les six paramètres
locaux aux centres ; résoudre sur les pixels originaux ; confirmer et propager ; ajouter
les nouvelles fenêtres ; réajuster. Arrêt sans gain ou après cinq tours. Aucun cas n'a
atteint la limite avec un gain encore en cours. Les observations déjà retenues ne sont
jamais remplacées par le champ lissé.

L'ajustement de ces petits systèmes est réalisé sur le processeur central. Le solveur
Gauss--Newton à composition inverse (IC-GN), utilisé pour la corrélation d'images
numériques (DIC), s'exécute sur WebGPU. Ce prototype expérimental n'est pas une chaîne
entièrement résidente sur le processeur graphique.

\newpage
\section*{2. Comparaison contrôlée et contribution réelle des gradients}
\fig{coverage.pdf}{75mm}
La phase initiale est identique pour toutes les variantes : résultats archivés de la
reprise affine avec rectification. La nouvelle étape essaie les propositions sur
\textbf{tous les centres éligibles}, contre une sous-grille dans l'expérience précédente.
La reprise affine de contrôle emploie exactement la même densité d'essais : elle isole
ce changement du choix de représentation.

\textbf{Contrôle hors ajustement.} Pour chacune des 5\,804 observations déjà retenues,
on retire à la fois son déplacement et son gradient, puis on les prédit avec les autres
observations. Il s'agit d'une cohérence avec les sorties du solveur, pas d'une vérité mécanique.

\begin{center}\begin{tabular}{lrrrr}
\toprule
 & \multicolumn{2}{c}{Écart de déplacement [pixels]} & \multicolumn{2}{c}{Écart de gradient, norme de Frobenius}\\
Champ & Médiane & Moyenne quadratique & Médiane & Moyenne quadratique\\\midrule
Affine &2,847&4,774&0,1875&0,2410\\
RBF déplacements &0,947&2,092&0,0900&0,1465\\
RBF déplacements et gradients &0,955&1,755&0,0778&0,1393\\\bottomrule
\end{tabular}\end{center}
Les gradients améliorent les erreurs quadratiques de prédiction, sans améliorer la
médiane du déplacement. Les deux RBF produisent néanmoins \textbf{exactement les mêmes
masques retenus} aux trois seuils sur les 191 paires. Il faut distinguer qualité du champ
proposé et nombre final de fenêtres acceptées.

\newpage
\section*{3. Où se trouvent les gains ? Exemple du niveau 191}
\fig{map_191.pdf}{125mm}
Au seuil 0,25, les RBF apportent deux fenêtres de plus que l'affine. Aux seuils 0,50 et
1,00, le choix de représentation ne change plus la couverture de ce cas. Le relâchement
ouvre des branches de propagation. Les fenêtres orange restantes au seuil 1,00 sont
65, centre $(47,5;79,5)$, et 164, centre $(143,5;175,5)$.

Les seuils photométrique (0,65), de déplacement (8 pixels), le maximum de 30 itérations,
la tolérance $10^{-3}$, l'interpolation bilinéaire et les supports restent inchangés.
Le seuil de gradient intervient dans la confirmation et dans la propagation, pas dans
l'acceptation numérique interne du solveur.

\newpage
\section*{4. Les images remappées expliquent l'intérêt du champ non affine}
\fig{remap_191.pdf}{70mm}
Les vues montrent $I_1(\mathbf x+\widehat{\mathbf u}(\mathbf x))$ dans le même rectangle
de référence de $440\times225$ pixels. Chaque vue repart de la cible originale. Ce
remappage est un \textbf{diagnostic visuel} : l'expérience numérique fournit directement
les six paramètres au solveur natif, sans résoudre sur une image RBF rééchantillonnée.
\fig{patch_191_102.pdf}{60mm}
Fenêtre 102 : l'affine donne une corrélation de 0,492 ; la première proposition RBF avec
gradients, 0,746 ; la dernière, 0,799 ; le solveur atteint 0,889. Le champ RBF rapproche
les structures avant résolution. Son rôle reste une initialisation ; la sortie retenue
provient du solveur, pas de l'interpolation du champ.

\newpage
\section*{5. Les deux échecs restants ne disparaissent pas avec le seuil de gradient}
\fig{patch_191_65.pdf}{60mm}
\textbf{Fenêtre 65.} La dernière sortie atteint une corrélation recalculée de 0,765,
mais garde le statut 0 : le critère de convergence n'est pas atteint dans la configuration
fixée. Une corrélation suffisante seule ne permet pas de la déclarer convergée. Le
rejet n'est plus celui du seuil de cohérence du gradient.
\fig{patch_191_164.pdf}{60mm}
\textbf{Fenêtre 164.} La proposition RBF reste très éloignée photométriquement de la
référence ; la dernière sortie n'atteint que 0,338 et le statut reste 0. Le nouveau
champ ne suffit donc pas à fournir toutes les graines. Ces images ne départagent pas
encore un mauvais bassin d'initialisation et une limite du modèle affine sur 32 pixels.

\newpage
\section*{6. Faut-il relâcher la cohérence du gradient ?}
\textbf{Oui pour l'expérimentation contrôlée :} le seuil absolu de 0,25 bloque ici de
nombreuses sorties que le solveur accepte. Au niveau 191, le seuil 0,50 récupère 13 fenêtres
au-delà de la phase affine et le seuil 1,00 en récupère 16. Une erreur de gradient est
sans dimension ; avec un étirement global voisin de 4,9, un seuil absolu n'a pas la même
signification relative que dans une transformation proche de l'identité.

\textbf{Pas comme nouveau défaut qualifié :} les 1\,719 configurations ne présentent
aucune fenêtre retenue hors image, sous 0,65 de corrélation finale ou à déterminant de
jacobienne non positif. Ces contrôles n'excluent pas tous les faux appariements. Aucun
changement des seuils par défaut du projet ou du solveur externe n'est effectué.
Le seuil 0,50 constitue un compromis à qualifier ; 1,00 reste une variante plus permissive.

\textbf{Ce que je retiens de la RBF contrainte par gradients.} Elle est préférable à
une affine comme représentation des observations non affines de cette plaque, et
la validation par retrait d'une observation confirme un intérêt des gradients. Mais
son avantage final sur les seules positions n'est pas démontré en nombre de fenêtres.
Les gaussiennes n'utilisent pas une distance suivant la matière : elles peuvent coupler
des observations de part et d'autre d'un trou. Ce modèle n'est pas qualifié pour transmettre
des graines à travers une fissure.

\textbf{Prochaine vérification ciblée :} qualifier les nouvelles acceptations à seuil
0,50 puis 1,00 sur des transformations connues à géométrie et hétérogénéité comparables,
avant de modifier une politique de production. Pour les deux fenêtres persistantes,
séparer un essai à support plus petit d'un essai disposant d'une meilleure initialisation,
sans modifier simultanément les critères de convergence.

\textbf{Coût et reproductibilité.} Les médianes de la seule nouvelle étape varient entre
0,044 et 0,120 seconde dans cette exécution, hors phase affine initiale. L'ordre fixe des
variantes, les caches et les écritures des traces interdisent d'en tirer un facteur
d'accélération général. Les scripts, paramètres, sorties par fenêtre, appels, contrôles
et empreintes sont conservés dans
\texttt{results/hudge\_rbf\_restart\_20261115/}. Dix tests ciblés et le contrôle de style
passent. Les dérivées analytiques sont vérifiées par différences finies, l'exactitude
affine et l'exclusion des fenêtres non retenues sont testées.
\end{document}
